\documentclass[a5paper,10pt]{article}

% https://github.com/InDevRus/filippov-solutions

% Нумерация страниц
\pagenumbering{gobble}

% Настройка полей
\usepackage{geometry}
\geometry{left=1cm}
\geometry{right=1cm}
\geometry{top=1cm}
\geometry{bottom=1.5cm}

% Вставка таблицы
\usepackage{array}
\usepackage{tabularx}

% Инструменты выравнивания формул
\usepackage{amsmath}
\usepackage{mathtools}

% Диакритика в математике
\usepackage{amsfonts}

% Вычисления размеров на ходу
\usepackage{calc}

% Удобные записи производных
\usepackage[italic=true]{derivative}

% Настройки сносок
\usepackage{enotez}

% Длинные знаки векторов
\usepackage{anyfontsize}
\usepackage[b]{esvect}

% Вставка картинок
\usepackage{graphicx}

% Вставка красной строки
\usepackage{indentfirst}
\setlength{\parindent}{1em}

% Условия в командах
\usepackage{xifthen}

% Луа-скрипты
\usepackage{luacode}

% Настройка команд отдельно в математическом и текстовом режимах
\usepackage{mathcommand}

% Для повторения LaTeX кода
\usepackage{multido}

% Конфигурация отступов формул
\delimitershortfall-1sp
\usepackage{mleftright}
\mleftright

% Растягивание объектов
\usepackage{relsize}

% Обтекание текстом
\usepackage{wrapfig}
\usepackage{subcaption}
\usepackage{floatflt}

% Поддержка ввода кириллицы
\usepackage[english, russian]{babel}

% Настройка корневой позиции для компилятора
\usepackage{import}
\providecommand*{\ProjectRootPath}{..}

\import{\ProjectRootPath/preambles/macros/}{theorems}

\import{\ProjectRootPath/preambles/fonts}{font_setup}

% Знак градуса
\usepackage{siunitx}

\import{\ProjectRootPath/preambles/macros/}{operators}
\import{\ProjectRootPath/preambles/macros/}{redefinitions}

\import{\ProjectRootPath/preambles/fonts}{letters_setup}
\import{\ProjectRootPath/preambles/fonts/adjustments}{custom_superscripts}

\import{\ProjectRootPath/preambles/custom}{formatting}
\import{\ProjectRootPath/preambles/custom}{frames}
\import{\ProjectRootPath/preambles/custom}{list_setup}

% TODO: перевести команды на современный стиль объявления

\begin{document}
    $\br{x - z} \parder{u} {x} + \br{\y - z} \parder{u} {y} + 2z \parder{u}{z} = 0$.

    $\dvfrac{\di x}{x - z} = \dvfrac{\di y}{y - z} = \dvfrac{\di z}{2z}$.

    \begin{enumerate}
        \item Из первой пропорции получаем производную пропорцию
        $\dvfrac{\di x - \di y}{x - y} = \dvfrac{\di z}{2z}$,
        из которой получаем первый интеграл
        $\ln\vbr{x - y} = \dfrac{1}{2}\ln\vbr{z} + A$.
        $\dfrac{\br{x - y}^{2}}{z} = A$.

        \item Из первых двух дробей можно получить ещё одну
        интегрируемую комбинацию путём линейного комбинирования числителей
        и знаменателей.
        Допустим, что мы комбинируем $\di x$, $a \di y$ и $b \di z$ 
        для некоторых $a$ и $b$. Тогда 
        $\dvfrac{\di x + a \di y + b \di z}{x + ay + \br{-1 - a + 2b} z}$
        будет являться полным дифференциалом, если и только если выполняется
        $2b - a - 1 = b$, то есть $b - a = 1$.
        Например, для $b = 2$, $a = 1$ получаем интегрируемую комбинацию
        $\dvfrac{\di x + \di y + 2 \di z}{x - z + y - z + 4z} = \dvfrac{\di z}{2z}$,
        откуда уже
        $\ln\vbr{x + y + 2z} = \dfrac{1}{2} \ln\vbr{z} + B$;
        $\dvfrac{\br{x + y + 2z}^2}{z} = B$.

        Стоит также отметить, что первая интегрируемая комбинация есть частный случай
        при $a = -1$, $b = 0$.
    \end{enumerate}

    Поскольку исходное уравнение является однородным, общее решение можно записать
    в форме
    $u = \f\br{\dfrac{\br{x - y}^{2}}{z}; \dvfrac{\br{x + y + 2z}^2}{z}}$.
\end{document}
